% ============================================================
%  Paper 12: PMNS Neutrino Mixing and the Solar-Isocurvature Cross-Sector Identity
%  Target: RevTeX 4-2 Compilation (SDGFT Series)
% ============================================================
\documentclass[amsmath,amssymb,aps,prd,twocolumn,superscriptaddress,nofootinbib]{revtex4-2}

\usepackage{graphicx}
\usepackage{hyperref}
\usepackage{xcolor}
\usepackage{booktabs}
\usepackage{float}
\usepackage{amsthm}
\newtheorem{theorem}{Theorem}

% ── SDGFT shared macros ──────────────────────────────────────
\usepackage{sdgft-macros}

% ── Document ─────────────────────────────────────────────────
\begin{document}

\preprint{SDGFT-2026-12}

\title{PMNS Neutrino Mixing and the Solar-Isocurvature Cross-Sector Identity}

\author{David A. Besemer}
\email{david.besemer@sdgft.org}
\affiliation{Independent Researcher}

\date{\today}

\begin{abstract}
We derive the PMNS neutrino mixing angles from the six-cone spatial partition of SDGFT. We establish the solar-isocurvature cross-sector identity, which couples the solar mixing angle sin2_theta_12 to the cosmic neutrino isocurvature fraction beta_iso = 1/36, yielding sin2_theta_12 = 11/36 ~ 0.3056. The atmospheric angle is shown to be sin2_theta_23 = 41/72 ~ 0.5694.
\end{abstract}

\maketitle

\section{Introduction}
Neutrino mixing angles are larger than quark mixing angles. We explain this difference through the direct projection of the spatial six-cone geometry.

\section{Theoretical Formulation}
Here we formulate the core mathematical relations governing the system:
\begin{equation}
  \sin^2\theta_{12} = \frac{1}{3} - \betaiso = \frac{1}{3} - \frac{1}{36} = \frac{11}{36} \approx 0.30556
\end{equation}
\begin{equation}
  \sin^2\theta_{23} = \frac{41}{72} \approx 0.56944
\end{equation}
\section{Cross-Sector Identity}
The solar angle is coupled to the early-universe isocurvature fraction, demonstrating that flavor physics and cosmology are intimately linked.

\section{Conclusion}
The PMNS angles are topologically locked, providing a testable prediction for neutrino oscillation experiments.



\begin{thebibliography}{99}
\bibitem{Goldstein_Paper01} D. A. Besemer, ``Axiomatic Foundations of SDGFT,'' SDGFT-2026-01 (in preparation).
\bibitem{Goldstein_Paper04} D. A. Besemer, ``Effective Spectral Dimension and the Banach Fixed-Point Theorem in SDGFT,'' SDGFT-2026-04.
\bibitem{Goldstein_Code} D. A. Besemer, \texttt{sdgft} Python package, \url{https://github.com/david-besemer/sdgft} (2026).
\end{thebibliography}

\end{document}
